
\begin{proposition}[Credited re-derivation: The non-holomorphic reduction theorem]\label{cdp:st:E2}
Let $Y$ be a compact connected complex threefold with $b_2(Y)=0$ and
let $g:Y\dashrightarrow C$ be a nonconstant meromorphic map to a
compact nonsingular curve which does not extend holomorphically.
For every holomorphic vector bundle $V$ on $Y$, there is a dense open
subset of each component of $\operatorname{Pic}(Y)$ on which
\[
 H^q(Y,V\otimes L)=0\quad(0\leq q\leq3).
\]
Consequently $\chi(Y,V\otimes L)=0$ for every $L\in\operatorname{Pic}(Y)$,
and $\int_Yc_3(Y)=0$.
\end{proposition}

The vanishing, Euler-index and third-Chern-number conclusions are \citet[Theorem~3.1, p.~9]{LRS2019}; its sphere consequence is Corollary~3.2 there. The extension from $\operatorname{Pic}^0$ to each component follows by tensoring the arbitrary vector bundle with a fixed representative of that component. We retain the complete graph proof below as a credited re-derivation, with the multiplicity filtration made explicit.

\begin{proof}
We give the graph proof, including the comparison needed at the
exceptional locus. The corresponding resolution comparison in the
printed proof of CDP20, Theorem~2.1, is justified separately in the
audit following \ref{cdp:st:E17}. The graph method is that of \citet[Theorem~3.1, p.~9]{LRS2019}; the multiplicity filtration below specifies its use on a full
scheme-theoretic fibre.

The curve $C$ is rational. Indeed resolve the map by successive blowups
with smooth centres. A map from a projective space to a positive-genus
curve is constant, so a holomorphic map from the last blowup to such a
curve is constant on its exceptional projective-space fibres. Its graph
then descends to the preceding smooth space: a proper bimeromorphic
graph with singleton fibres over a normal space is an isomorphism.
Descending at each stage would extend $g$ on $Y$, contrary to its
definition. Take the Stein factor of a smooth resolution of the map
and replace $g$ by the induced meromorphic map to that factor. It
is still non-holomorphic, since otherwise the original map would
be holomorphic. The same descent argument makes its target
$\mathbb P^1$. Its general fibre is irreducible and reduced;
its image in the graph below need not be smooth.

Let $Z\subset Y\times\mathbb P^1$ be the closure of its graph and write
\[
 \pi:Z\to Y,\qquad h:Z\to\mathbb P^1
\]
for the holomorphic projections. Locally the map is $[a:b]$, where
$a,b$ are holomorphic germs with no common prime factor. The graph has
equation $av-bu=0$ in coordinates $[u:v]$ on $\mathbb P^1$; this
hypersurface is integral. On this local product the exact hypersurface
sequence is
\[
 0\to\mathcal O(-1)\xrightarrow{\ av-bu\ }\mathcal O
 \to\mathcal O_Z\to0,
\]
with the harmless line bundle from $Y$ included in the first term when
the local equations are changed. The cohomology of $\mathcal O$ and
$\mathcal O(-1)$ on $\mathbb P^1$ gives
$R\pi_*\mathcal O_Z=\mathcal O_Y$. Moreover $R^q\pi_*\mathcal F=0$
for $q\geq2$ and every coherent $\mathcal F$, since every fibre of
$\pi$ has dimension at most one. For every closed subspace $T\subset Z$,
the ideal sequence and these two facts imply
$R^1\pi_*\mathcal O_T=0$. These are statements about this graph
projection, with their explicitly given local equations.

The two sections defining the meromorphic map are sections of a line
bundle $P$ on $Y$ with no common divisorial component. Their pullbacks
generate an invertible ideal on $Z$. It defines an effective Cartier
divisor $\Sigma$ with
\begin{equation}\label{cdp:eq:E-graph-pencil}
 h^*\mathcal O_{\mathbb P^1}(1)=\pi^*P\otimes\mathcal O_Z(-\Sigma).
\end{equation}
The local ideal is generated by $a/u$ or $b/v$ on the corresponding
graph chart, proving both the asserted invertibility and the displayed
comparison. Every component of $\Sigma$ dominates $\mathbb P^1$:
its support lies over the codimension-two common-zero locus of $a,b$,
and the fibre over each such point is the entire projective line.
In particular $\Sigma$ meets every fibre. A fibre $Z_c$ is identified
by $\pi$ with the Cartier divisor $(a v_c-bu_c=0)$ in $Y$, including
its multiplicities. It is therefore Cohen--Macaulay and Gorenstein.

We recall the elementary vanishing used below. On a compact reduced
irreducible space $D$, for a locally free sheaf $W$ and a nonzero
effective Cartier divisor $A$, the descending finite-dimensional
spaces $H^0(D,W(-kA))$ stabilize. A section in their stable value
vanishes to every order along $A$, hence vanishes by the Krull
intersection theorem and the identity theorem. The stable value is
zero. On a Cartier divisor with multiplicities, successively remove
one prime component at a time. The corresponding exact sequences
have graded factors $\mathcal O_D(-D_{\rm previous})$ on reduced prime
components; tensoring by an invertible sheaf and a vector bundle
preserves exactness. Thus it suffices to apply the preceding vanishing
to the finitely many reduced-component factors.

We now choose $M\in\operatorname{Pic}(Y)$ and $k>0$ such that
\begin{equation}\label{cdp:eq:E-all-fibres}
 H^0\bigl(Z_c,\pi^*(V^*\otimes M^{-1})\otimes
       \omega_{Z_c}\otimes\mathcal O_{Z_c}(-k\Sigma)\bigr)=0
 \quad\text{for every }c.
\end{equation}
Here is a simultaneous construction. There are only finitely many
reducible or nonreduced fibres. Form the graph of their reduced prime
components, joining two which intersect in $Y$. Give a component
weight equal to its shortest distance from a component meeting
$\Sigma$ on $Z$. Every weight is finite because the fibres are
connected and meet $\Sigma$. Let $D_i$ be the sum, with coefficient
one, of the components of weight $i$. The inclusion $D_i\hookrightarrow
Y$ is a Cartier divisor since $Y$ is smooth.

Put $M=\mathcal O_Y(\sum_{i<d}n_iD_i)$, where $d$ is the largest
weight. Choose $n_{d-1}$ first, then $n_{d-2}$, and continue down to
$n_0$. On each component of weight $i>0$, $D_{i-1}$ restricts to a
nonzero effective Cartier divisor; $D_j$ with $j<i-1$ is disjoint
from it. After the higher-index coefficients have been fixed, the
preceding vanishing permits a choice of $n_{i-1}$ which annihilates
the sections on every multiplicity-filtration factor belonging to
weight $i$. Later choices do not affect those factors. The remaining
weight-zero factors are annihilated by increasing $k$, because
$\Sigma$ restricts to a nonzero effective Cartier divisor on each.
This proves \eqref{cdp:eq:E-all-fibres} on the finitely many exceptional
fibres.

For clarity, uniformity on the other fibres does not presume local
freeness of any higher direct image. First choose a power of
$-\Sigma$ annihilating sections on one general irreducible reduced fibre. Serre
duality and top-degree base change show that
$R^2h_*(\pi^*(V\otimes P^k))$ then has finite support. Increasing $k$
can only decrease that support, because multiplication by the
nonzero divisor cutting out $\Sigma$ injects the dual section spaces.
At each of its finitely many irreducible fibres the same vanishing
lemma, with its multiplicity filtration, kills the remaining group.
The effective twist $M$ cannot introduce a section on another fibre:
each of its prime divisors restricts there to an effective Cartier
divisor. This completes the simultaneous construction.

Restriction of \eqref{cdp:eq:E-graph-pencil} to $Z_c$ gives the actual line
isomorphism $\pi^*P|_{Z_c}=\mathcal O_{Z_c}(\Sigma)$. Serre duality
on the Gorenstein surface $Z_c$ turns \eqref{cdp:eq:E-all-fibres} into
\[
 H^2(Z_c,\pi^*(V\otimes M\otimes P^k)|_{Z_c})=0.
\]
Top-degree proper base change and Nakayama's lemma now give
$R^2h_*\pi^*(V\otimes M\otimes P^k)=0$. Serre vanishing on
$\mathbb P^1$ supplies $t$ for which
\[
 H^2\bigl(Z,\pi^*(V\otimes M\otimes P^k)
                       \otimes h^*\mathcal O(t)\bigr)=0.
\]
By \eqref{cdp:eq:E-graph-pencil}, this is
$H^2(Z,\pi^*(V\otimes M\otimes P^{k+t})(-t\Sigma))=0$.
The sheaf on $t\Sigma$ has zero second cohomology: its pushforward to
$Y$ is supported on the at-most-one-dimensional exceptional image,
and its higher direct images vanish by the graph calculation above
and the projection formula. The divisor exact sequence therefore
gives
$H^2(Z,\pi^*(V\otimes M\otimes P^{k+t}))=0$. Finally
$R\pi_*\mathcal O_Z=\mathcal O_Y$ and the projection formula give
an isomorphism with $H^2(Y,V\otimes M\otimes P^{k+t})$.

Since $H^2(Y,\mathbb Z)$ is finite, all exponents in this construction
can be chosen divisible by its exponent. The resulting twisting line
lies in $\operatorname{Pic}^0(Y)$. The same construction applied to
$V\otimes L_*$ reaches any fixed component $L_*\operatorname{Pic}^0(Y)$.
For degree zero use the vanishing lemma directly with any nonzero
effective fibre divisor on $Y$ and large negative multiples.
For degrees three and one use Serre duality, respectively with these
degree-zero and degree-two arguments applied to $V^*\otimes K_Y$.
Semicontinuity in the holomorphic exponential family of line bundles
over $H^1(Y,\mathcal O_Y)$ gives proper analytic jumping loci; their
finite union has dense open complement. The exponential family is
constructed by exponentiating a basis of additive Cech one-cocycles,
so it includes every line in $\operatorname{Pic}^0(Y)$ and the
semicontinuity assertion needs no compactness of that group.

All line bundles have rationally zero first Chern class. Hence
Riemann--Roch makes $\chi(V\otimes L)$ independent of $L$, and the
just-proved vanishing makes it zero. For $V=T_Y$, the terms involving
$c_1(Y)$ vanish in the index, leaving
$\chi(T_Y)=\tfrac12\int_Yc_3(Y)$. This proves the last assertion.
\end{proof}


